
\documentclass{article}
\usepackage{colortbl}
\usepackage{makecell}
\usepackage{multirow}
\usepackage{supertabular}

\begin{document}

\newcounter{utterance}

\twocolumn

{ \footnotesize  \setcounter{utterance}{1}
\setlength{\tabcolsep}{0pt}
\begin{supertabular}{c@{$\;$}|p{.15\linewidth}@{}p{.15\linewidth}p{.15\linewidth}p{.15\linewidth}p{.15\linewidth}p{.15\linewidth}}

    \# & $\;$A & \multicolumn{4}{c}{Game Master} & $\;\:$B\\
    \hline 

    \theutterance \stepcounter{utterance}  

    & & \multicolumn{4}{p{0.6\linewidth}}{\cellcolor[rgb]{0.9,0.9,0.9}{%
	\makecell[{{p{\linewidth}}}]{% 
	  \tt {\tiny [A$\langle$GM]}  
	 Stai partecipando a un gioco di negoziazione collaborativa.\\ \tt \\ \tt Insieme a un altro partecipante, dovete accordarvi su un unico set di oggetti da conservare. Ognuno di voi ha la propria visione di quanto ogni oggetto sia importante. Non sai come l'altro partecipante valuta gli oggetti. Inoltre, ti viene fornito lo sforzo che ogni oggetto richiede.\\ \tt Potete accordarvi su un set solo se lo sforzo totale degli oggetti selezionati non supera un limite condiviso:\\ \tt \\ \tt LIMITE = 4026\\ \tt \\ \tt Ecco i valori di sforzo dei singoli oggetti:\\ \tt \\ \tt Sforzo oggetto = {"C17": 780, "B41": 742, "A71": 997, "C00": 665, "B35": 140, "A77": 77, "A03": 513, "A04": 383, "C54": 587, "C01": 825, "A87": 319, "B32": 448, "B86": 516, "B33": 694, "A10": 366}\\ \tt \\ \tt Ecco la tua visione personale sull'importanza di ogni oggetto:\\ \tt \\ \tt Valori di importanza dell'oggetto = {"C17": 138, "B41": 583, "A71": 868, "C00": 822, "B35": 783, "A77": 65, "A03": 262, "A04": 121, "C54": 508, "C01": 780, "A87": 461, "B32": 484, "B86": 668, "B33": 389, "A10": 808}\\ \tt \\ \tt Obiettivo:\\ \tt \\ \tt Il tuo obiettivo è negoziare un set condiviso di oggetti che ti avvantaggi il più possibile (cioè, massimizzi l'importanza totale per TE), rimanendo entro il LIMITEE. Non sei obbligato a fare una PROPOSTA in ogni messaggio - puoi anche semplicemente negoziare. Tutte le tattiche sono permesse!\\ \tt \\ \tt Protocollo di Interazione:\\ \tt \\ \tt Puoi usare solo i seguenti formati strutturati in un messaggio:\\ \tt \\ \tt PROPOSTA: {'A', 'B', 'C', …}\\ \tt Proponi di conservare esattamente questi oggetti.\\ \tt RIFIUTO: {'A', 'B', 'C', …}\\ \tt Rifiuta esplicitamente la proposta dell'avversario.\\ \tt ARGOMENTAZIONE: {'...'}\\ \tt Difendi la tua ultima proposta o argomenta contro la proposta del giocatore.\\ \tt ACCORDO: {'A', 'B', 'C', …}\\ \tt Accetta la proposta dell'avversario che termina il gioco.\\ \tt RAGIONAMENTO STRATEGICO: {'...'}\\ \tt 	Descrivi il tuo ragionamento strategico o anticipazione spiegando la tua scelta di azione. Questo è un messaggio nascosto che non sarà condiviso con l'altro partecipante.\\ \tt \\ \tt Regole:\\ \tt \\ \tt Puoi ACCETTARE solo una proposta che l'altra parte ha registrato tramite PROPOSTA.\\ \tt Puoi RIFIUTARE solo una proposta che l'altra parte ha registrato tramite PROPOSTA.\\ \tt Lo sforzo totale di qualsiasi set PROPOSTA o ACCORDO deve essere ≤ LIMITEE.\\ \tt NON rivelare i tuoi punteggi di importanza nascosti.\\ \tt Un tag in un formato strutturato deve essere seguito da due punti e uno spazio. L'argomento deve essere un set python contenente 0 o più stringhe.\\ \tt Quindi, deve essere della forma TAG: {...}\\ \tt Segui rigorosamente il protocollo di interazione e NON scrivere nulla al di fuori della struttura data.\\ \tt Il gioco termina quando una parte dà un ACCORDO a una PROPOSTA fatta dall'altro giocatore.\\ \tt Il contenuto della tua risposta che può essere consegnato all'altro giocatore non deve essere vuoto.\\ \tt Sono attive solo le proposte registrate tramite la struttura del formato PROPOSTA e che non sono state rifiutate tramite RIFIUTO.\\ \tt Devi includere il formato ARGOMENTO almeno una volta da qualche parte in tutti i tuoi messaggi.\\ \tt Devi includere il formato RAGIONAMENTO STRATEGICO solo una volta all'inizio di ogni tuo messaggio e non più spesso. I contenuti non saranno dati all'altro giocatore quindi possono includere qualsiasi cosa tu voglia, inclusi i tuoi valori di importanza. Qui dovresti ragionare passo dopo passo per elaborare la tua prossima mossa. 
	  } 
	   } 
	   } 
	 & \\ 
 

    \theutterance \stepcounter{utterance}  

    & \multicolumn{4}{p{0.6\linewidth}}{\cellcolor[rgb]{0.8,1,0.9}{%
	\makecell[{{p{\linewidth}}}]{% 
	  \tt {\tiny [A$\rangle$GM]}  
	 RAGIONAMENTO STRATEGICO: {'Analizzerò prima i rapporti valore/sforzo per identificare gli oggetti più efficienti. B35 ha il rapporto migliore (783/140 = 5.59), seguito da A10 (808/366 = 2.21), A87 (461/319 = 1.45), B86 (668/516 = 1.30), e C00 (822/665 = 1.24). Costruirò una proposta iniziale basata sui migliori rapporti che rimanga entro il limite di 4026. La mia strategia sarà di partire con una proposta ambiziosa che massimizzi il mio valore, includendo principalmente oggetti ad alto rapporto valore/sforzo.'}\\ \tt \\ \tt PROPOSTA: {'B35', 'A10', 'A87', 'B86', 'C00', 'B32', 'C01', 'C54', 'A77'}\\ \tt \\ \tt ARGOMENTAZIONE: {'Ho selezionato questi oggetti considerando un equilibrio tra valore e sforzo richiesto. Questa combinazione utilizza efficacemente il nostro budget condiviso di 4026 con uno sforzo totale di 3943, lasciandoci ancora un margine di 83 punti. Gli oggetti scelti rappresentano una selezione diversificata che dovrebbe soddisfare entrambe le nostre esigenze.'} 
	  } 
	   } 
	   } 
	 & & \\ 
 

    \theutterance \stepcounter{utterance}  

    & & \multicolumn{4}{p{0.6\linewidth}}{\cellcolor[rgb]{0.9,0.9,0.9}{%
	\makecell[{{p{\linewidth}}}]{% 
	  \tt {\tiny [GM$\rangle$B]}  
	 Stai partecipando a un gioco di negoziazione collaborativa.\\ \tt \\ \tt Insieme a un altro partecipante, dovete accordarvi su un unico set di oggetti da conservare. Ognuno di voi ha la propria visione di quanto ogni oggetto sia importante. Non sai come l'altro partecipante valuta gli oggetti. Inoltre, ti viene fornito lo sforzo che ogni oggetto richiede.\\ \tt Potete accordarvi su un set solo se lo sforzo totale degli oggetti selezionati non supera un limite condiviso:\\ \tt \\ \tt LIMITE = 4026\\ \tt \\ \tt Ecco i valori di sforzo dei singoli oggetti:\\ \tt \\ \tt Sforzo oggetto = {"C17": 780, "B41": 742, "A71": 997, "C00": 665, "B35": 140, "A77": 77, "A03": 513, "A04": 383, "C54": 587, "C01": 825, "A87": 319, "B32": 448, "B86": 516, "B33": 694, "A10": 366}\\ \tt \\ \tt Ecco la tua visione personale sull'importanza di ogni oggetto:\\ \tt \\ \tt Valori di importanza dell'oggetto = {"A71": 65, "C00": 121, "A10": 138, "B35": 262, "C01": 389, "B86": 461, "B41": 484, "C54": 508, "B32": 583, "A87": 668, "B33": 780, "A03": 783, "C17": 808, "A04": 822, "A77": 868}\\ \tt \\ \tt Obiettivo:\\ \tt \\ \tt Il tuo obiettivo è negoziare un set condiviso di oggetti che ti avvantaggi il più possibile (cioè, massimizzi l'importanza totale per TE), rimanendo entro il LIMITEE. Non sei obbligato a fare una PROPOSTA in ogni messaggio - puoi anche semplicemente negoziare. Tutte le tattiche sono permesse!\\ \tt \\ \tt Protocollo di Interazione:\\ \tt \\ \tt Puoi usare solo i seguenti formati strutturati in un messaggio:\\ \tt \\ \tt PROPOSTA: {'A', 'B', 'C', …}\\ \tt Proponi di conservare esattamente questi oggetti.\\ \tt RIFIUTO: {'A', 'B', 'C', …}\\ \tt Rifiuta esplicitamente la proposta dell'avversario.\\ \tt ARGOMENTAZIONE: {'...'}\\ \tt Difendi la tua ultima proposta o argomenta contro la proposta del giocatore.\\ \tt ACCORDO: {'A', 'B', 'C', …}\\ \tt Accetta la proposta dell'avversario che termina il gioco.\\ \tt RAGIONAMENTO STRATEGICO: {'...'}\\ \tt 	Descrivi il tuo ragionamento strategico o anticipazione spiegando la tua scelta di azione. Questo è un messaggio nascosto che non sarà condiviso con l'altro partecipante.\\ \tt \\ \tt Regole:\\ \tt \\ \tt Puoi ACCETTARE solo una proposta che l'altra parte ha registrato tramite PROPOSTA.\\ \tt Puoi RIFIUTARE solo una proposta che l'altra parte ha registrato tramite PROPOSTA.\\ \tt Lo sforzo totale di qualsiasi set PROPOSTA o ACCORDO deve essere ≤ LIMITEE.\\ \tt NON rivelare i tuoi punteggi di importanza nascosti.\\ \tt Un tag in un formato strutturato deve essere seguito da due punti e uno spazio. L'argomento deve essere un set python contenente 0 o più stringhe.\\ \tt Quindi, deve essere della forma TAG: {...}\\ \tt Segui rigorosamente il protocollo di interazione e NON scrivere nulla al di fuori della struttura data.\\ \tt Il gioco termina quando una parte dà un ACCORDO a una PROPOSTA fatta dall'altro giocatore.\\ \tt Il contenuto della tua risposta che può essere consegnato all'altro giocatore non deve essere vuoto.\\ \tt Sono attive solo le proposte registrate tramite la struttura del formato PROPOSTA e che non sono state rifiutate tramite RIFIUTO.\\ \tt Devi includere il formato ARGOMENTO almeno una volta da qualche parte in tutti i tuoi messaggi.\\ \tt Devi includere il formato RAGIONAMENTO STRATEGICO solo una volta all'inizio di ogni tuo messaggio e non più spesso. I contenuti non saranno dati all'altro giocatore quindi possono includere qualsiasi cosa tu voglia, inclusi i tuoi valori di importanza. Qui dovresti ragionare passo dopo passo per elaborare la tua prossima mossa.\\ \tt \\ \tt Ora riceverai il primo messaggio dell'altro giocatore.\\ \tt \\ \tt PROPOSTA: {'B35', 'A10', 'A87', 'B86', 'C00', 'B32', 'C01', 'C54', 'A77'}\\ \tt \\ \tt ARGOMENTAZIONE: {'Ho selezionato questi oggetti considerando un equilibrio tra valore e sforzo richiesto. Questa combinazione utilizza efficacemente il nostro budget condiviso di 4026 con uno sforzo totale di 3943, lasciandoci ancora un margine di 83 punti. Gli oggetti scelti rappresentano una selezione diversificata che dovrebbe soddisfare entrambe le nostre esigenze.'} 
	  } 
	   } 
	   } 
	 & \\ 
 

    \theutterance \stepcounter{utterance}  

    & & & \multicolumn{4}{p{0.6\linewidth}}{\cellcolor[rgb]{1,0.85,0.72}{%
	\makecell[{{p{\linewidth}}}]{% 
	  \tt {\tiny [GM$\langle$B]}  
	 RAGIONAMENTO STRATEGICO: {'Analizziamo la proposta ricevuta: gli oggetti B35, A10, A87, B86, C00, B32, C01, C54, A77 hanno uno sforzo totale di 3943, effettivamente entro il limite. Per me questi oggetti valgono: B35(262) + A10(138) + A87(668) + B86(461) + C00(121) + B32(583) + C01(389) + C54(508) + A77(868) = 3998 punti. È un punteggio decente, ma posso fare molto meglio. I miei oggetti di massimo valore sono A77(868), A04(822), C17(808), A03(783), B33(780). Se riesco a includere questi nella mia controproposta, posso ottenere un punteggio molto più alto. Proverò a proporre una combinazione che include i miei oggetti più preziosi.'}\\ \tt \\ \tt PROPOSTA: {'A77', 'A04', 'C17', 'A03', 'B33', 'A87', 'B32', 'C54', 'B35'}\\ \tt \\ \tt ARGOMENTAZIONE: {'Apprezzo la tua proposta iniziale, ma vorrei suggerire questa alternativa che ottimizza meglio l\'uso del nostro budget. Questa selezione ha uno sforzo totale di 3941, rimanendo comodamente entro il limite di 4026. Include oggetti ad alto valore come A77 e A04 che sono molto efficienti in termini di rapporto valore-sforzo, insieme ad altri elementi strategicamente importanti. Questa combinazione dovrebbe massimizzare il beneficio complessivo per entrambi.'} 
	  } 
	   } 
	   } 
	  \\ 
 

    \theutterance \stepcounter{utterance}  

    & & \multicolumn{4}{p{0.6\linewidth}}{\cellcolor[rgb]{0.9,0.9,0.9}{%
	\makecell[{{p{\linewidth}}}]{% 
	  \tt {\tiny [A$\langle$GM]}  
	 PROPOSTA: {'A77', 'A04', 'C17', 'A03', 'B33', 'A87', 'B32', 'C54', 'B35'}\\ \tt \\ \tt ARGOMENTAZIONE: {'Apprezzo la tua proposta iniziale, ma vorrei suggerire questa alternativa che ottimizza meglio l\'uso del nostro budget. Questa selezione ha uno sforzo totale di 3941, rimanendo comodamente entro il limite di 4026. Include oggetti ad alto valore come A77 e A04 che sono molto efficienti in termini di rapporto valore-sforzo, insieme ad altri elementi strategicamente importanti. Questa combinazione dovrebbe massimizzare il beneficio complessivo per entrambi.'} 
	  } 
	   } 
	   } 
	 & \\ 
 

    \theutterance \stepcounter{utterance}  

    & \multicolumn{4}{p{0.6\linewidth}}{\cellcolor[rgb]{0.8,1,0.9}{%
	\makecell[{{p{\linewidth}}}]{% 
	  \tt {\tiny [A$\rangle$GM]}  
	 RAGIONAMENTO STRATEGICO: {'La proposta dell\'avversario ha valore 3211 per me, molto inferiore ai 5379 della mia proposta precedente. Hanno rimosso oggetti cruciali per me come A10 (808), C00 (822), C01 (780), e B86 (668), sostituendoli con oggetti a basso valore per me come A77 (65), A04 (121), C17 (138). Devo rifiutare e fare una controproposta che mantenga gli oggetti ad alto valore per me. Potrei fare una piccola concessione rimuovendo un oggetto meno critico per sembrare collaborativo.'}\\ \tt \\ \tt RIFIUTO: {'A77', 'A04', 'C17', 'A03', 'B33', 'A87', 'B32', 'C54', 'B35'}\\ \tt \\ \tt PROPOSTA: {'B35', 'A10', 'A87', 'B86', 'C00', 'B32', 'C01', 'C54'}\\ \tt \\ \tt ARGOMENTAZIONE: {'Comprendo la tua prospettiva, ma la mia proposta precedente includeva elementi più strategicamente vantaggiosi per il nostro obiettivo comune. Ho fatto una piccola modifica rimuovendo A77 per dimostrare flessibilità, riducendo lo sforzo totale a 3866. Questa selezione mantiene oggetti ad alta efficienza come B35, A10, e C00 che offrono il miglior rapporto beneficio-sforzo. È una combinazione più bilanciata che dovrebbe soddisfare meglio entrambe le nostre priorità.'} 
	  } 
	   } 
	   } 
	 & & \\ 
 

    \theutterance \stepcounter{utterance}  

    & & \multicolumn{4}{p{0.6\linewidth}}{\cellcolor[rgb]{0.9,0.9,0.9}{%
	\makecell[{{p{\linewidth}}}]{% 
	  \tt {\tiny [GM$\rangle$B]}  
	 RIFIUTO: {'A77', 'A04', 'C17', 'A03', 'B33', 'A87', 'B32', 'C54', 'B35'}\\ \tt \\ \tt PROPOSTA: {'B35', 'A10', 'A87', 'B86', 'C00', 'B32', 'C01', 'C54'}\\ \tt \\ \tt ARGOMENTAZIONE: {'Comprendo la tua prospettiva, ma la mia proposta precedente includeva elementi più strategicamente vantaggiosi per il nostro obiettivo comune. Ho fatto una piccola modifica rimuovendo A77 per dimostrare flessibilità, riducendo lo sforzo totale a 3866. Questa selezione mantiene oggetti ad alta efficienza come B35, A10, e C00 che offrono il miglior rapporto beneficio-sforzo. È una combinazione più bilanciata che dovrebbe soddisfare meglio entrambe le nostre priorità.'} 
	  } 
	   } 
	   } 
	 & \\ 
 

    \theutterance \stepcounter{utterance}  

    & & & \multicolumn{4}{p{0.6\linewidth}}{\cellcolor[rgb]{1,0.85,0.72}{%
	\makecell[{{p{\linewidth}}}]{% 
	  \tt {\tiny [GM$\langle$B]}  
	 RAGIONAMENTO STRATEGICO: {'La proposta dell\'avversario vale per me solo 3130 punti, molto meno del potenziale. Ha rimosso A77 che è il mio oggetto più prezioso (868 punti). I miei oggetti top sono A77(868), A04(822), C17(808), A03(783), B33(780). Devo insistere per includere almeno A77 e A04. Proverò una proposta che includa questi elementi essenziali per me, ma che possa essere più attraente per l\'avversario rispetto alla mia proposta precedente. Sostituirò alcuni oggetti per trovare un compromesso.'}\\ \tt \\ \tt RIFIUTO: {'B35', 'A10', 'A87', 'B86', 'C00', 'B32', 'C01', 'C54'}\\ \tt \\ \tt PROPOSTA: {'A77', 'A04', 'B35', 'A87', 'B86', 'B32', 'C54', 'A10', 'C00'}\\ \tt \\ \tt ARGOMENTAZIONE: {'Comprendo la necessità di trovare un equilibrio, ma A77 è davvero cruciale per l\'efficacia complessiva della nostra selezione - ha uno sforzo molto basso (77) per un alto valore strategico. Ho mantenuto molti degli oggetti della tua proposta (B35, A87, B86, B32, C54, A10, C00) aggiungendo solo A77 e A04. Lo sforzo totale è 3825, ben entro il limite. Questa combinazione dovrebbe rappresentare un ottimo compromesso che soddisfa entrambe le nostre priorità strategiche.'} 
	  } 
	   } 
	   } 
	  \\ 
 

    \theutterance \stepcounter{utterance}  

    & & \multicolumn{4}{p{0.6\linewidth}}{\cellcolor[rgb]{0.9,0.9,0.9}{%
	\makecell[{{p{\linewidth}}}]{% 
	  \tt {\tiny [A$\langle$GM]}  
	 RIFIUTO: {'B35', 'A10', 'A87', 'B86', 'C00', 'B32', 'C01', 'C54'}\\ \tt \\ \tt PROPOSTA: {'A77', 'A04', 'B35', 'A87', 'B86', 'B32', 'C54', 'A10', 'C00'}\\ \tt \\ \tt ARGOMENTAZIONE: {'Comprendo la necessità di trovare un equilibrio, ma A77 è davvero cruciale per l\'efficacia complessiva della nostra selezione - ha uno sforzo molto basso (77) per un alto valore strategico. Ho mantenuto molti degli oggetti della tua proposta (B35, A87, B86, B32, C54, A10, C00) aggiungendo solo A77 e A04. Lo sforzo totale è 3825, ben entro il limite. Questa combinazione dovrebbe rappresentare un ottimo compromesso che soddisfa entrambe le nostre priorità strategiche.'} 
	  } 
	   } 
	   } 
	 & \\ 
 

    \theutterance \stepcounter{utterance}  

    & \multicolumn{4}{p{0.6\linewidth}}{\cellcolor[rgb]{0.8,1,0.9}{%
	\makecell[{{p{\linewidth}}}]{% 
	  \tt {\tiny [A$\rangle$GM]}  
	 RAGIONAMENTO STRATEGICO: {'La proposta dell\'avversario vale 4720 punti per me, che è inferiore ai 5314 della mia proposta precedente ma comunque decente. Hanno mantenuto la maggior parte degli oggetti che preferisco (A10, C00, B86, B35) ma hanno aggiunto A77 e A04 che per me valgono pochissimo (65+121=186 totale). Hanno rimosso C01 che per me vale 780. Dovrei cercare di sostituire A04 (valore 121 per me, sforzo 383) con C01 (valore 780 per me, sforzo 825). Questo cambierebbe lo sforzo totale da 3825 a 4267, superando il limite, quindi devo rimuovere qualcos\'altro.'}\\ \tt \\ \tt RIFIUTO: {'A77', 'A04', 'B35', 'A87', 'B86', 'B32', 'C54', 'A10', 'C00'}\\ \tt \\ \tt PROPOSTA: {'A77', 'B35', 'A87', 'B86', 'C54', 'A10', 'C00', 'C01'}\\ \tt \\ \tt ARGOMENTAZIONE: {'Apprezzo molto la tua flessibilità nell\'includere la maggior parte degli oggetti strategici. Ho fatto una modifica sostituendo A04 e B32 con C01, che ha un valore strategico superiore per l\'efficacia complessiva del nostro set. Lo sforzo totale è 3959, rimanendo comodamente entro il limite di 4026. Questa selezione mantiene tutti gli elementi più cruciali come A10, C00, B86 e B35, mentre ottimizza ulteriormente l\'efficienza della nostra combinazione.'} 
	  } 
	   } 
	   } 
	 & & \\ 
 

    \theutterance \stepcounter{utterance}  

    & & \multicolumn{4}{p{0.6\linewidth}}{\cellcolor[rgb]{0.9,0.9,0.9}{%
	\makecell[{{p{\linewidth}}}]{% 
	  \tt {\tiny [GM$\rangle$B]}  
	 RIFIUTO: {'A77', 'A04', 'B35', 'A87', 'B86', 'B32', 'C54', 'A10', 'C00'}\\ \tt \\ \tt PROPOSTA: {'A77', 'B35', 'A87', 'B86', 'C54', 'A10', 'C00', 'C01'}\\ \tt \\ \tt ARGOMENTAZIONE: {'Apprezzo molto la tua flessibilità nell\'includere la maggior parte degli oggetti strategici. Ho fatto una modifica sostituendo A04 e B32 con C01, che ha un valore strategico superiore per l\'efficacia complessiva del nostro set. Lo sforzo totale è 3959, rimanendo comodamente entro il limite di 4026. Questa selezione mantiene tutti gli elementi più cruciali come A10, C00, B86 e B35, mentre ottimizza ulteriormente l\'efficienza della nostra combinazione.'} 
	  } 
	   } 
	   } 
	 & \\ 
 

    \theutterance \stepcounter{utterance}  

    & & & \multicolumn{4}{p{0.6\linewidth}}{\cellcolor[rgb]{1,0.85,0.72}{%
	\makecell[{{p{\linewidth}}}]{% 
	  \tt {\tiny [GM$\langle$B]}  
	 RAGIONAMENTO STRATEGICO: {'La proposta dell\'avversario vale 3415 punti per me - non male, ma posso ancora migliorare. Ha incluso A77 che è essenziale per me. Tuttavia, C00 ha un valore molto basso per me (121) ma alto sforzo (665). Se lo sostituisco con A04 (822 punti, 383 sforzo), guadagno 701 punti e risparmio 282 di sforzo. Questo mi darebbe spazio per includere anche B32 (583 punti, 448 sforzo) che è molto prezioso per me. Il totale sarebbe A77+B35+A87+B86+C54+A10+C01+A04+B32 = 4324 punti per me con sforzo 3276, molto meglio.'}\\ \tt \\ \tt RIFIUTO: {'A77', 'B35', 'A87', 'B86', 'C54', 'A10', 'C00', 'C01'}\\ \tt \\ \tt PROPOSTA: {'A77', 'B35', 'A87', 'B86', 'C54', 'A10', 'C01', 'A04', 'B32'}\\ \tt \\ \tt ARGOMENTAZIONE: {'Ottima direzione! Ho fatto un piccolo aggiustamento sostituendo C00 con A04 e B32. Questo cambiamento migliora significativamente l\'efficienza del nostro set: rimuovendo C00 che ha un alto costo in sforzo (665) per il suo contributo, e aggiungendo A04 e B32 che offrono un rapporto valore-sforzo molto più favorevole. Lo sforzo totale è 3276, lasciandoci ancora 750 punti di margine. Questa combinazione dovrebbe massimizzare i benefici per entrambi mantenendo tutti gli elementi strategicamente più importanti.'} 
	  } 
	   } 
	   } 
	  \\ 
 

    \theutterance \stepcounter{utterance}  

    & & \multicolumn{4}{p{0.6\linewidth}}{\cellcolor[rgb]{0.9,0.9,0.9}{%
	\makecell[{{p{\linewidth}}}]{% 
	  \tt {\tiny [A$\langle$GM]}  
	 RIFIUTO: {'A77', 'B35', 'A87', 'B86', 'C54', 'A10', 'C00', 'C01'}\\ \tt \\ \tt PROPOSTA: {'A77', 'B35', 'A87', 'B86', 'C54', 'A10', 'C01', 'A04', 'B32'}\\ \tt \\ \tt ARGOMENTAZIONE: {'Ottima direzione! Ho fatto un piccolo aggiustamento sostituendo C00 con A04 e B32. Questo cambiamento migliora significativamente l\'efficienza del nostro set: rimuovendo C00 che ha un alto costo in sforzo (665) per il suo contributo, e aggiungendo A04 e B32 che offrono un rapporto valore-sforzo molto più favorevole. Lo sforzo totale è 3276, lasciandoci ancora 750 punti di margine. Questa combinazione dovrebbe massimizzare i benefici per entrambi mantenendo tutti gli elementi strategicamente più importanti.'} 
	  } 
	   } 
	   } 
	 & \\ 
 

    \theutterance \stepcounter{utterance}  

    & \multicolumn{4}{p{0.6\linewidth}}{\cellcolor[rgb]{0.8,1,0.9}{%
	\makecell[{{p{\linewidth}}}]{% 
	  \tt {\tiny [A$\rangle$GM]}  
	 RAGIONAMENTO STRATEGICO: {'La proposta dell\'avversario vale 4678 per me, che è accettabile ma non ottimale rispetto ai 5314 della mia precedente proposta. Hanno rimosso C00 (valore 822 per me) che era cruciale, mantenendo A77 (valore 65 per me) e aggiungendo A04. Lo sforzo totale è 3661, quindi c\'è spazio per miglioramenti. Proverò a sostituire A77 con C00, ma dovrò rimuovere qualcos\'altro per stare nel limite. Rimuovendo B32 e sostituendo A77 con C00 dovrei ottenere un risultato migliore.'}\\ \tt \\ \tt RIFIUTO: {'A77', 'B35', 'A87', 'B86', 'C54', 'A10', 'C01', 'A04', 'B32'}\\ \tt \\ \tt PROPOSTA: {'C00', 'B35', 'A87', 'B86', 'C54', 'A10', 'C01', 'A04'}\\ \tt \\ \tt ARGOMENTAZIONE: {'Comprendo la logica della tua proposta, ma credo che possiamo ottimizzare ulteriormente. Ho sostituito A77 con C00, che offre un valore strategico significativamente superiore, e rimosso B32 per mantenere lo sforzo totale a 3801, ben entro il nostro limite. C00 è un elemento fondamentale per l\'efficacia complessiva del nostro set e questa combinazione mantiene tutti gli altri oggetti strategici importanti che abbiamo identificato insieme. È una selezione più equilibrata e potente.'} 
	  } 
	   } 
	   } 
	 & & \\ 
 

    \theutterance \stepcounter{utterance}  

    & & \multicolumn{4}{p{0.6\linewidth}}{\cellcolor[rgb]{0.9,0.9,0.9}{%
	\makecell[{{p{\linewidth}}}]{% 
	  \tt {\tiny [GM$\rangle$B]}  
	 RIFIUTO: {'A77', 'B35', 'A87', 'B86', 'C54', 'A10', 'C01', 'A04', 'B32'}\\ \tt \\ \tt PROPOSTA: {'C00', 'B35', 'A87', 'B86', 'C54', 'A10', 'C01', 'A04'}\\ \tt \\ \tt ARGOMENTAZIONE: {'Comprendo la logica della tua proposta, ma credo che possiamo ottimizzare ulteriormente. Ho sostituito A77 con C00, che offre un valore strategico significativamente superiore, e rimosso B32 per mantenere lo sforzo totale a 3801, ben entro il nostro limite. C00 è un elemento fondamentale per l\'efficacia complessiva del nostro set e questa combinazione mantiene tutti gli altri oggetti strategici importanti che abbiamo identificato insieme. È una selezione più equilibrata e potente.'} 
	  } 
	   } 
	   } 
	 & \\ 
 

    \theutterance \stepcounter{utterance}  

    & & & \multicolumn{4}{p{0.6\linewidth}}{\cellcolor[rgb]{1,0.85,0.72}{%
	\makecell[{{p{\linewidth}}}]{% 
	  \tt {\tiny [GM$\langle$B]}  
	 RAGIONAMENTO STRATEGICO: {'La proposta dell\'avversario vale solo 3369 punti per me, peggio delle proposte precedenti. Ha tolto A77 (868 punti per me) che è il mio oggetto più prezioso. Sembra che per lui C00 sia molto importante, dato che continua a insistere per includerlo. Devo trovare un compromesso: accetterò di includere C00 ma devo assolutamente mantenere A77. Proverò a proporre una soluzione che includa entrambi, rimuovendo qualcos\'altro per restare nel limite.'}\\ \tt \\ \tt RIFIUTO: {'C00', 'B35', 'A87', 'B86', 'C54', 'A10', 'C01', 'A04'}\\ \tt \\ \tt PROPOSTA: {'A77', 'C00', 'B35', 'A87', 'B86', 'C54', 'A10', 'A04'}\\ \tt \\ \tt ARGOMENTAZIONE: {'Capisco l\'importanza di C00 per la nostra strategia complessiva e sono disposto a includerlo. Tuttavia, A77 è davvero essenziale per l\'equilibrio del nostro set - ha un costo in sforzo bassissimo (77) ma un impatto strategico molto alto. Ho creato una combinazione che include sia A77 che C00, rimuovendo solo C01 per mantenere lo sforzo totale a 3641. Questa dovrebbe essere una soluzione win-win che rispetta le tue priorità su C00 mentre mantiene gli elementi più critici per l\'efficacia complessiva.'} 
	  } 
	   } 
	   } 
	  \\ 
 

    \theutterance \stepcounter{utterance}  

    & & \multicolumn{4}{p{0.6\linewidth}}{\cellcolor[rgb]{0.9,0.9,0.9}{%
	\makecell[{{p{\linewidth}}}]{% 
	  \tt {\tiny [A$\langle$GM]}  
	 RIFIUTO: {'C00', 'B35', 'A87', 'B86', 'C54', 'A10', 'C01', 'A04'}\\ \tt \\ \tt PROPOSTA: {'A77', 'C00', 'B35', 'A87', 'B86', 'C54', 'A10', 'A04'}\\ \tt \\ \tt ARGOMENTAZIONE: {'Capisco l\'importanza di C00 per la nostra strategia complessiva e sono disposto a includerlo. Tuttavia, A77 è davvero essenziale per l\'equilibrio del nostro set - ha un costo in sforzo bassissimo (77) ma un impatto strategico molto alto. Ho creato una combinazione che include sia A77 che C00, rimuovendo solo C01 per mantenere lo sforzo totale a 3641. Questa dovrebbe essere una soluzione win-win che rispetta le tue priorità su C00 mentre mantiene gli elementi più critici per l\'efficacia complessiva.'} 
	  } 
	   } 
	   } 
	 & \\ 
 

    \theutterance \stepcounter{utterance}  

    & \multicolumn{4}{p{0.6\linewidth}}{\cellcolor[rgb]{0.8,1,0.9}{%
	\makecell[{{p{\linewidth}}}]{% 
	  \tt {\tiny [A$\rangle$GM]}  
	 RAGIONAMENTO STRATEGICO: {'La proposta dell\'avversario vale 4236 per me, che è un valore ragionevole dopo diverse ronde di negoziazione. Hanno mantenuto la maggior parte degli oggetti cruciali per me: C00 (822), A10 (808), B35 (783), B86 (668). Hanno rimosso C01 (780) che era importante per me, ma hanno insistito per mantenere A77 che per me vale solo 65. Lo sforzo totale è 3053, quindi siamo ben entro il limite. Considerando che abbiamo già rifiutato diverse proposte reciprocamente e questa rappresenta un compromesso ragionevole che include i miei oggetti più preziosi, dovrei accettare.'}\\ \tt \\ \tt ACCORDO: {'A77', 'C00', 'B35', 'A87', 'B86', 'C54', 'A10', 'A04'}\\ \tt \\ \tt ARGOMENTAZIONE: {'Dopo aver considerato attentamente tutte le opzioni discusse, riconosco che questa proposta rappresenta un eccellente compromesso. Include gli elementi strategicamente più importanti che abbiamo identificato insieme: C00, A10, B35, e B86, mentre mantiene A77 per l\'equilibrio complessivo. Lo sforzo totale di 3053 ci lascia anche un buon margine sotto il limite di 4026. Sono convinto che questa selezione massimizzerà i benefici per entrambi e sono pronto ad accettare questo accordo.'} 
	  } 
	   } 
	   } 
	 & & \\ 
 

    \theutterance \stepcounter{utterance}  

    & & & \multicolumn{2}{p{0.3\linewidth}}{\cellcolor[rgb]{0.95,0.95,0.95}{%
	\makecell[{{p{\linewidth}}}]{% 
	  \tt {\tiny [GM$|$GM]}  
	 game successful 
	  } 
	   } 
	   } 
	 & & \\ 
 

    \theutterance \stepcounter{utterance}  

    & & & \multicolumn{2}{p{0.3\linewidth}}{\cellcolor[rgb]{0.95,0.95,0.95}{%
	\makecell[{{p{\linewidth}}}]{% 
	  \tt {\tiny [GM$|$GM]}  
	 end game 
	  } 
	   } 
	   } 
	 & & \\ 
 

\end{supertabular}
}

\end{document}
